\section{Sensitivitaetsanalyse}
\label{sec:sensitivitaet}

% VERANTWORTLICH: ...
% Aufgabenpunkt 6.

\subsection{Vorgehen}

% One-at-a-time: ein Parameter wird variiert, die uebrigen bleiben fest.
% Grenzen dieses Vorgehens benennen: Wechselwirkungen zwischen Parametern
% werden damit nicht erfasst.
% Bewertungsgroessen festlegen, z.B. Maximaltemperatur und Energieertrag.

\subsection{Ergebnisse}

\subsection{Einordnung}

% Guter Diskussionspunkt: Tuncel et al. finden bei Stundenwerten praktisch
% keinen Einfluss der Waermekapazitaet und schliessen daraus, dass eine
% stationaere Rechnung genuegen wuerde. Herteleer et al. argumentieren fuer
% Sekunden- bis Minutenaufloesung ueber die thermische Zeitkonstante
% tau = R_eq * C_eq von einigen hundert Sekunden.
% Mit den 10-Minuten-Daten liegt diese Arbeit dazwischen. Zeigen, auf
% welcher Seite das eigene Ergebnis landet, und daraus ableiten, welche
% Parameter ueberhaupt sorgfaeltig recherchiert werden muessen.
